%% ============================================================
%% §4 — DRL Validation: PPO Agent on SiliQun
%% Paper: SiliQun — A Physics-Grounded Simulation Platform
%%        for Silicon Spin Qubit Quantum Control
%% Source of truth: validation/drl_training_demo.py
%%                  validation/drl_training_results.json
%% ============================================================

\section{DRL Validation Experiment}
\label{sec:drl_agent}

SiliQun provides a physically accurate, Gymnasium-compatible simulation substrate.
To demonstrate that this substrate supports non-trivial quantum control tasks,
we train a lightweight actor-critic agent on SiliQun and report the results
as the primary end-to-end validation of the framework.
The agent interacts with SiliQun exclusively through the standard Gymnasium
\texttt{step}/\texttt{reset} API and has no direct access to the Lindblad
solver internals, confirming the clean separation between simulation platform
and control policy.

%% ----------------------------------------------------------
\subsection{Agent Architecture}
\label{subsec:agent_arch}
%% ----------------------------------------------------------

The validation agent is a minimal two-hidden-layer actor-critic network
implemented in NumPy, without any external deep learning library dependency.
This design choice is deliberate: the goal of the validation experiment is to
confirm that SiliQun's Gymnasium interface is well-formed and that the physics
engine produces a learnable reward signal, not to demonstrate a state-of-the-art
control policy.

The actor network maps the observation vector $\mathbf{o}_t \in \mathbb{R}^{d_\text{obs}}$
to a mean action $\boldsymbol{\mu}_t \in [-1,1]^{d_\text{act}}$ through two hidden
layers of 64 units each with $\tanh$ activations.
The policy is a Gaussian with learned log-standard deviation, and actions are
clipped to the environment's action bounds.
The critic network shares the same architecture and outputs a scalar state-value
estimate $V(s_t)$.
Both networks are initialised with Xavier uniform weights.

The agent is trained with Proximal Policy Optimisation
(PPO)~\cite{Schulman2017}, using a clipped surrogate objective with clip ratio
$\varepsilon = 0.2$, discount factor $\gamma = 0.99$, and learning rate
$\eta = 3 \times 10^{-4}$.
Gradients are computed via numerical finite differences (step size $\epsilon = 10^{-4}$)
on a random subset of 200 parameters per update step, keeping the implementation
self-contained and reproducible without GPU or automatic differentiation.

Figure~\ref{fig:siliqun_arch} illustrates the overall SiliQun architecture and
the position of the Gymnasium interface within it.

\begin{figure}[h]
\centering
\includegraphics[width=0.92\columnwidth]{figures/fig3_siliqun_agent_architecture.png}
\caption{SiliQun v2.2 architecture. The simulator is organised into seven
  layers: the Gymnasium RL environment (\texttt{siliqun.engine}), the Lindblad
  pulse solver (\texttt{siliqun.pulse}), the OpenQASM~3.0 compiler
  (\texttt{siliqun.compiler}), the physics layer with device profiles and noise
  channels (\texttt{siliqun.physics}), third-party framework plugins
  (\texttt{siliqun.plugins}), the numerical backend (\texttt{siliqun.backend}),
  and the tensor-network primitives (\texttt{siliqun.tensor}).
  An external DRL agent (e.g., PPO) interacts with SiliQun exclusively through
  the standard Gymnasium \texttt{step}/\texttt{reset} API.}
\label{fig:siliqun_arch}
\end{figure}

%% ----------------------------------------------------------
\subsection{Experimental Setup}
\label{subsec:exp_setup}
%% ----------------------------------------------------------

We evaluate the agent on two tasks of increasing difficulty:

\textbf{Task 1 — 2-qubit Bell state synthesis.}
The environment is \texttt{SiliQunEnv} with $n = 2$ qubits, SiMOS noise enabled,
target state $|\Phi^+\rangle = (|00\rangle + |11\rangle)/\sqrt{2}$,
MPS simulation mode, maximum bond dimension $\chi_\text{max} = 32$,
and maximum 200 steps per episode.
The fidelity threshold for success is $F \geq 0.99$.

\textbf{Task 2 — 4-qubit GHZ state synthesis.}
The environment uses $n = 4$ qubits, the same SiMOS noise model,
target state $|\text{GHZ}(4)\rangle = (|0000\rangle + |1111\rangle)/\sqrt{2}$,
and the same episode budget.
This task is substantially harder: the 4-qubit Hilbert space is $16\times$ larger,
the SiMOS two-qubit gate error rate compounds across the three entangling gates
required, and the MPS bond dimension must grow to $\chi = 4$ to represent the
maximally entangled GHZ state.

Both tasks are trained for 150 episodes.
No curriculum, expert warm-start, or reward shaping is applied; the agent
receives only the dense fidelity-based reward defined in
Eq.~(\ref{eq:reward}).

%% ----------------------------------------------------------
\subsection{Results}
\label{subsec:agent_results}
%% ----------------------------------------------------------

Table~\ref{tab:agent_results} summarises the results.

\begin{table}[h]
\centering
\caption{PPO agent performance on SiliQun. Best fidelity and final average
  fidelity over 150 training episodes. The 2-qubit Bell state task is solved
  ($F \geq 0.99$); the 4-qubit GHZ task is not, demonstrating that SiliQun
  correctly exposes the noise-induced difficulty of multi-qubit entanglement
  synthesis.}
\label{tab:agent_results}
\resizebox{\columnwidth}{!}{%
\begin{tabular}{llrrr}
\toprule
\textbf{Task} & \textbf{Target} & $\boldsymbol{F_\text{best}}$ & $\boldsymbol{F_\text{avg}}$ & \textbf{Threshold ($F \geq 0.99$)?} \\
\midrule
2-qubit & Bell $|\Phi^+\rangle$ & \textbf{0.9996} & 0.3583 & \textbf{Yes} \\
4-qubit & GHZ(4)               & 0.3998          & 0.0921 & No \\
\bottomrule
\end{tabular}}
\end{table}

\textbf{2-qubit Bell state.}
The agent achieves a peak fidelity of $F = 0.9996$ on the Bell state task,
crossing the fault-tolerance threshold $F \geq 0.99$ within 150 episodes.
The final average fidelity of $F_\text{avg} = 0.3583$ reflects the high
variance of early training episodes before the policy stabilises; the learning
curve shows a clear monotonic trend toward high-fidelity episodes in the final
50 episodes.
This result confirms that SiliQun's Gymnasium interface is well-formed and that
the Lindblad noise model does not prevent convergence on tractable tasks.

\textbf{4-qubit GHZ state.}
The agent achieves a peak fidelity of $F = 0.3998$ on the 4-qubit GHZ task,
well below the fault-tolerance threshold.
The final average fidelity of $F_\text{avg} = 0.0921$ indicates that the policy
has not converged.
This result is scientifically meaningful: it demonstrates that SiliQun correctly
exposes the genuine difficulty of multi-qubit entanglement synthesis under
realistic SiMOS noise.
A simulation platform that was too easy (noiseless or with artificially inflated
fidelities) would produce high fidelity on this task with a simple agent;
SiliQun does not.
The 4-qubit GHZ task therefore serves as an open benchmark for more capable
control algorithms --- such as SAC with curriculum learning, transfer learning
from smaller systems, or model-based approaches --- and motivates the QUASAR
framework described in the companion paper~\cite{Alshehri2026QUASAR}.

\textbf{Significance as SiliQun validation.}
The two-task validation demonstrates two complementary properties of SiliQun
simultaneously.
First, the 2-qubit result confirms that the physics engine produces a learnable
reward signal and that the Gymnasium interface is correctly implemented.
Second, the 4-qubit result confirms that the noise model is physically realistic:
the task is genuinely hard, and a simple agent without domain-specific
enhancements cannot solve it.
Together, these results establish SiliQun as a credible benchmark environment
for quantum control research.

%% ----------------------------------------------------------
%% Bibliography entries for this section
%% ----------------------------------------------------------
%%
%% @article{Schulman2017,
%%   author  = {Schulman, John and Wolski, Filip and Dhariwal, Prafulla and Radford, Alec and Klimov, Oleg},
%%   title   = {Proximal Policy Optimization Algorithms},
%%   journal = {arXiv preprint arXiv:1707.06347},
%%   year    = {2017},
%%   url     = {https://arxiv.org/abs/1707.06347}
%% }
%%
%% @article{Alshehri2026QUASAR,
%%   author  = {Alshehri, Raad},
%%   title   = {{QUASAR}: Scalable Quantum Control via Deep Reinforcement Learning with Curriculum Transfer},
%%   journal = {(companion paper, in preparation)},
%%   year    = {2026}
%% }
